% Preface
\section*{Preface}\label{sec:preface}
\addcontentsline{toc}{section}{Preface}

\subsection*{Who this book is for}

This tutorial is written for \textbf{applied scientists and data scientists} with some mathematical background --- you have taken linear algebra and calculus, you use \texttt{lm()}, \texttt{LinearRegression()}\allowbreak, or \texttt{np.linalg.lstsq} on a daily basis, and yet you feel a thin pane of glass between ``knowing how to use'' and ``understanding'': why is the slope exactly the correlation coefficient? Why squared loss rather than absolute value? How does a model with zero training error in the $d > n$ regime manage to generalize at all? What exactly is the relationship between kernel methods and my linear regression?

If these questions sound familiar, this book is for you. It does not teach you how to call packages (that is what documentation is for); it derives every formula from scratch and hands you the \textbf{assumption list} behind every theorem --- because in statistics the assumptions are the substance, and the formulas are merely their shadow. The only prerequisites are undergraduate calculus and linear algebra plus basic probability; every claim can be verified by hand, and complete solutions to all exercises are collected in the appendix.

\subsection*{Origins and influences}

This book was written \textbf{on a whim} --- one evening, the author decided to write down, in one sitting, the ``why do we even learn linear regression'' argument that had been circulating in his head for years, and ended up going chapter by chapter all the way to Chapter~8. But the whim stands on the accumulated influence of a body of work read over a long time. The seven scholars who influenced it most are: \textbf{Francis Bach} (Inria), whose textbook \emph{Learning Theory from First Principles} (2024)\cite{bach2024} shares this book's temperament of earning every result from first principles --- the kernel-methods perspective of Chapter~\ref{sec:kernel} and the optimization--statistics interlocking of Chapters~\ref{sec:gd}--\ref{sec:gd-underdet} both bear its mark; \textbf{Ryan Tibshirani} (Stanford University), who turned generalized lasso, trend filtering, and conformal prediction into methodology that is both computable and auditable --- precisely the spirit of the estimate-to-inference route in Chapter~\ref{sec:ggm}; \textbf{Jingfeng Wu} (UC Berkeley, Simons Institute), whose work showing that ``gradient descent dominates ridge regression''\cite{wu2025} is the statistical frontier of the iteration-as-regularization story of Chapter~\ref{sec:gd} through Section~\ref{sec:implicit}; \textbf{Arthur Jacot} (NYU Courant Institute), whose neural tangent kernel\cite{jacot2018} welded ``linear regression $\to$ kernel regression $\to$ deep networks'' into a single road --- the content of Section~\ref{sec:ntk}; \textbf{Jianqing Fan} (Princeton University), whose high-dimensional statistical framework\cite{fan2020} --- sparse estimation, variable screening, large covariance matrices --- is the source of the from-estimation-to-inference roadmap of Chapter~\ref{sec:ggm}; \textbf{David Donoho} (Stanford University), founder of compressed sensing, who proved that $\ell_1$ minimization exploits sparsity optimally and that underdetermined systems become solvable under sparsity assumptions\cite{donoho2006} --- Section~\ref{sec:sparsity}'s ``sparsity is a computable assumption'' is the statistical expression of exactly this idea; \textbf{Emmanuel Cand\`es} (Stanford University), who co-created the exact-recovery theory of compressed sensing with Donoho and Tao, and who, through the knockoffs framework, gave complex-model variable selection reproducible false discovery rate (FDR) control\cite{barber2015} --- the modern next station after Section~\ref{sec:debiased}'s ``from point estimates to testable inference.''

\subsection*{How it was written, and an honest disclaimer}

\textbf{This book was written with substantial AI assistance.} Specifically, the text was produced by Moonshot AI's Kimi K2.8 model running in the \textbf{Kimi Work desktop} agent mode --- it wrote the LaTeX source directly, compiled the PDF, rendered and inspected pages, and proofread chapter by chapter; the human expert (Haoyi Xiong) supplied the outline, the chapter content, and his own thinking --- in particular his multi-year understanding of linear regression and the specific points he had personally worked on (such as the exact solution of gradient descent in the underdetermined case, the kernel-methods view under latent projections, and the precision-matrix--conditional-independence connection). One sentence of division of labor: \textit{the human decides ``what to ask'' and ``what to believe''; the machine does the ``writing'' and the ``checking.''}

Readers should therefore keep three things in mind:

\begin{enumerate}[label=(\alph*)]
  \item \textbf{The field of linear regression is vast}, and this book does not and cannot cover it all: generalized linear models, robust regression, mixed-effects models, time series, the Bayesian route, online learning, federated settings --- each deserves its own deep-dive. This book covers only the main line that the authors consider impossible to skip honestly.
  \item \textbf{The theory here does not aim for professional rigor.} Every formula is correct and derivable, but the proofs are plainly not textbook-complete (regularity conditions for asymptotic results are uniformly omitted). \textit{Readers who work on theory should skip this book}; beginners or those who have never seen the material systematically may read it as a popularization with derivations.
  \item \textbf{The historical and bibliographic sections were checked with care} (principal figures, years, and sources were verified online and marked in the bibliography), but if any error remains, the responsibility lies with the human authors, and corrections are welcome.
\end{enumerate}

Fit the straight line first, audit the assumptions, disclose the weaknesses honestly --- this book applies its own recipe to itself.

\vspace{1em}
\hfill Haoyi Xiong (熊昊一), in Beijing

\hfill Lingkai Kong (孔令凯), in Seattle
